\section{Rational Hankel sums and a finite polylogarithm alphabet}
\label{rha:section}

The rational-pole input of the preceding section resums a differential
normal form. Rational poles also control a different, finite-cutoff
construction: Hankel determinants of moments. Here the size of a
determinant can grow while its generating function remains in one fixed
polylogarithm alphabet. This provides an exact framework for experiments
with harmonic determinants, complementary to the retained rank problems
of Section~\ref{opf:section}. The confluent-pole report
\cite{rha:reports} supplies the theorem reconciled below. Harmonic degree,
polylogarithmic weight and arithmetic independence remain distinct.

Let $K$ be a field of characteristic zero embedded in $\mathbb C$, and let
\begin{equation}\label{rha:weight}
 R(x)=P(x)+\sum_{a\in A}\sum_{\ell=1}^{s_a}
                 \frac{c_{a,\ell}}{(x+a)^\ell},
 \qquad A\subset\mathbb Z_{\ge0}\text{ finite},\quad c_{a,s_a}\ne0,
\end{equation}
with $P\in K[x]$ and $c_{a,\ell}\in K$. Put
\begin{equation}\label{rha:det}
 D_{r,M}=\det\left(\sum_{n=1}^{M}n^{i+j}R(n)\right)_{0\le i,j<r},
 \qquad G_r(z)=\sum_{M\ge1}D_{r,M}z^M,\quad r\ge1.
\end{equation}
Finite Cauchy--Binet, or direct expansion of two Vandermonde
determinants, gives
\begin{equation}\label{rha:ensemble}
 D_{r,M}=\sum_{1\le n_1<\cdots<n_r\le M}
       \prod_{i<j}(n_j-n_i)^2\prod_{j=1}^{r}R(n_j).
\end{equation}
There is no $r!$ factor. In particular $D_{r,M}=0$ for $M<r$.
The pole restriction avoids every summation node, and the terms grow at
most polynomially in $M$. Hence $G_r$ is holomorphic for $|z|<1$.

\subsection{Confluent poles and the weight budget}

Write $h_{a,p}(M)=\sum_{n=1}^{M}(n+a)^{-p}$ and first regard these
as independent formal variables of weight $p$. For a pole of order $s_a$,
let $V_a$ have entries
\[
 V_a(i,p)=\binom ip(-a)^{i-p}\quad(0\le p\le i),\qquad
 V_a(i,p)=0\quad(p>i),\quad 0\le p<s_a,
\]
where $0^0=1$, and define
\[
 C_a(p,q)=\sum_{\ell>p+q}^{s_a}c_{a,\ell}h_{a,\ell-p-q}.
\]
Expansion of $x^i$ and $x^j$ around $-a$ proves the exact factorization
\begin{equation}\label{rha:factor}
 \boldsymbol\mu_r(M)=B_r(M)+\sum_{a\in A}V_aC_aV_a^{\mathsf T},
                  \qquad B_r(M)\in\operatorname{Mat}_r(K[M]).
\end{equation}
Indeed the terms with $p+q<\ell$ give $C_a$; all others have polynomial
Faulhaber sums. The transpose is ordinary transpose, including for
complex residues. A nonzero entry of $C_a$ has weight at most $s_a-p-q$.

\begin{theorem}[Confluent pole allocation]\label{rha:capacity}
Every monomial in the formal determinant has harmonic weight at most
\begin{equation}\label{rha:budget}
 W_r=\max_{\substack{0\le k_a\le s_a\\\sum_a k_a\le r}}
              \sum_a k_a(s_a-k_a+1).
\end{equation}
Equivalently, $W_r$ is the sum of the $r$ largest entries in the multiset
of positive integers $s_a,s_a-2,s_a-4,\ldots$ from all poles, or all
entries if fewer than $r$ are available. Thus
\[
 W_r\le W_\infty=\sum_a\left\lfloor\frac{(s_a+1)^2}{4}\right\rfloor.
\]
For no poles both budgets are zero.
\end{theorem}
\begin{proof}
Concatenate the $V_a$ and use block diagonal $C$ in \eqref{rha:factor}.
Multilinearity selects $k\le r$ columns from $VCV^{\mathsf T}$ and
the remainder from $B_r$. Cauchy--Binet then places all harmonic
dependence in minors $C[I,J]$ of size $k$. A nonzero block-diagonal
minor selects equally many indices $k_a$ from the row and column sets
of each pole. Its local weight is bounded by
\[
 k_as_a-\sum_{p\in I_a}p-\sum_{q\in J_a}q
 \le k_as_a-k_a(k_a-1)=k_a(s_a-k_a+1).
\]
The $V$ and $B_r$ factors have weight zero. Maximize over the possible
selections to obtain \eqref{rha:budget}. The successive increments of
$k(s-k+1)$ are $s,s-2,s-4,\ldots$. They decrease in each block, so
selecting the largest positive increments respects every block's prefix
constraint. Zero increments can be omitted; negative ones cannot improve
the maximum. Their sum proves the saturation formula.
\end{proof}

Simple poles have $W_r=\min(r,|A|)$. A single pole of orders
$1,2,3,4,5,6,7$ has stabilized weights $1,2,4,6,9,12,16$.
Thus neither ordinary harmonic degree nor pole order alone supplies
the required polylogarithmic weight. For example the two-particle
triple-pole determinant is $H_M^{(3)}H_M^{(1)}-(H_M^{(2)})^2$,
of harmonic degree two and weight four.

\subsection{Cutoff corrections and anchored rational integration}

The exact integer-shift identity is
\begin{equation}\label{rha:shift}
 h_{a,p}(M)=H_M^{(p)}+\sum_{b=1}^{a}(M+b)^{-p}-H_a^{(p)}.
\end{equation}
Assign cost $p$ to a retained harmonic factor or an inverse-cutoff
factor of order $p$, and zero to the rational constant. Expanding
\eqref{rha:shift} cannot increase total cost. Partial fractions and
polynomial division therefore write $D_{r,M}$ as a finite sum of
\begin{equation}\label{rha:cutoff}
 p(M)\prod_jH_M^{(b_j)},\qquad
 \frac{c}{(M+h)^e}\prod_jH_M^{(b_j)},
\end{equation}
where $p\in K[M]$, $c\in K$, $1\le h\le\max A$, and respectively
$\sum_jb_j\le W_r$ or $e+\sum_jb_j\le W_r$.
Each partial-fraction pole order is at most the sum of the preceding
denominator orders. The empty harmonic product is one; there are no
cutoff denominators if $A$ is empty or $A=\{0\}$.

Use decreasing-index finite sums and multiple polylogarithms:
\[
 H_M(b_1,\ldots,b_v)=\sum_{M\ge n_1>\cdots>n_v\ge1}
             \frac{1}{n_1^{b_1}\cdots n_v^{b_v}},\qquad
 \operatorname{Li}_{b_1,\ldots,b_v}(z)=
       \sum_{n_1>\cdots>n_v\ge1}
             \frac{z^{n_1}}{n_1^{b_1}\cdots n_v^{b_v}}.
\]
All indices are positive. Partitioning a product according to equal
summation indices gives the positive quasi-shuffle: for instance
$H_M(a)H_M(b)=H_M(a,b)+H_M(b,a)+H_M(a+b)$.
The recursion preserves total weight and integer multiplicities.
Summing the largest-index geometric tail proves
\begin{equation}\label{rha:hgen}
 \sum_{M\ge1}H_M(\mathbf b)z^M
       =\frac{\operatorname{Li}_{\mathbf b}(z)}{1-z}
                    \quad(\mathbf b\ne\emptyset).
\end{equation}
The empty word instead has generator $z/(1-z)$.

For clarity, the integration step needed by the compiler is proved here.
Let $I_\emptyset=1$, with kernels $f_0=1/z$ and $f_1=1/(1-z)$,
and $I_{a\nu}(z)=\int_0^z f_a(t)I_\nu(t)\,dt$ for nonempty words
ending in $1$. Then
$\operatorname{Li}_{b_1,\ldots,b_v}=I_{0^{b_1-1}1\cdots0^{b_v-1}1}$.

\begin{lemma}[Anchored two-letter Hermite reduction]\label{rha:hermite}
Suppose $F$ is a rational linear combination of $1$ and words ending
in $1$, of lengths at most $w$. Its rational coefficients have poles
only at $0,1$. If $F$ is holomorphic at zero, its zero-anchored
primitive has the same form with lengths at most $w+1$.
An exact terminating algorithm constructs it.
\end{lemma}
\begin{proof}
Partial fractions give $R=Q'+\alpha/z+\beta/(1-z)$ with rational $Q$
having only those poles. For a nonempty word $\omega=a\nu$,
\[
 \int R I_\omega\,dz=Q I_\omega+\alpha I_{0\omega}
       +\beta I_{1\omega}-\int Q f_a I_\nu\,dz.
\]
The remaining word is shorter, so recursion terminates. At the empty
word the primitive is $Q+\alpha\log z-\beta\log(1-z)$.
After combining terms, the result is a meromorphic germ plus a constant
multiple of $\log z$; all other words are analytic at zero.
Its derivative is holomorphic by hypothesis, forcing both the principal
part and the $\log z$ coefficient to vanish. Only finitely many Taylor
coefficients are needed to compute the constant term, because the
rational factors have finite pole order. Subtract that constant.
This also proves the assertion for $F/z$ when $F(0)=0$.
\end{proof}

This is the two-letter specialization of classical hyperlogarithm
integration; quasi-shuffle and rational-integration attribution are
retained in \cite{rha:reports}. Define $\theta=z\,d/dz$ and
\begin{equation}\label{rha:inverse}
 \mathcal J_hF(z)=z^{-h}\int_0^z t^{h-1}F(t)\,dt,
                \qquad h\in\mathbb Z_{\ge0}.
\end{equation}
For $h=0$ require $F(0)=0$; for $h>0$ the apparent pole is removable.
Coefficientwise, $\mathcal J_h\sum_{M\ge1}a_Mz^M
=\sum_{M\ge1}a_Mz^M/(M+h)$. The operator $\theta$ preserves the
weight bound, whereas $\mathcal J_h$ raises it by at most one.

\subsection{An all-rank finite alphabet}

\begin{theorem}[Finite rational-pole compiler]\label{rha:compiler}
For every weight \eqref{rha:weight} and rank $r\ge1$, $G_r(z)$ is a
finite $K(z)$-linear combination of $1$ and positive-index multiple
polylogarithms of weight at most $W_r$. The prefactors have poles only
at $0,1$. The identity holds throughout $|z|<1$, with removable values
at zero taken collectively.
\end{theorem}
\begin{proof}
Apply the confluent expansion, the shift rule and cutoff partial
fractions \eqref{rha:cutoff}. Quasi-shuffle every harmonic product and
use \eqref{rha:hgen}. Multiplication by $p(M)$ becomes $p(\theta)$;
division by $(M+h)^e$ becomes $\mathcal J_h^e$. Lemma~\ref{rha:hermite}
removes every integral using exact rational operations and a recursion
of strictly decreasing word length. A term of harmonic weight $b$ and
inverse-cutoff order $e$ produces weight at most $b+e\le W_r$.
All constructions are finite and introduce no new letter. Their
coefficient identities establish equality of Taylor series, while
analyticity extends equality to the full disk.
\end{proof}

\begin{corollary}[Every denominator of degree at most two]
\label{rha:quadratic}
If the reduced denominator of $R$ has degree at most two and roots at
nonpositive integers, then, at every rank and with any polynomial part,
\begin{equation}\label{rha:quadratic-form}
 G_r(z)=A_r(z)+B_r(z)L+C_r(z)L^2+E_r(z)\operatorname{Li}_2(z),
                         \qquad L=-\log(1-z),
\end{equation}
with prefactors in $K(z)$ having poles only at $0,1$.
\end{corollary}
\begin{proof}
The nonempty pole profiles are $(1)$, $(1,1)$ and $(2)$, with
$W_\infty\le2$. The positive-index polylogarithms of weight at most
two are $\operatorname{Li}_1=L$, $\operatorname{Li}_2$, and
$\operatorname{Li}_{1,1}=L^2/2$. Include the rational empty word.
\end{proof}

Separated integer poles and repeated double poles therefore have the
same uniform dilogarithm bound as adjacent simple poles. This is a
constructive upper bound, not a minimal functional-weight theorem.

\begin{corollary}[Euler primitives and differential equations]
\label{rha:differential}
For $e\ge0$, the generator $\sum_{M\ge1}D_{r,M}z^M/M^e$ has weight
at most $W_r+e$. All ranks for fixed $R$ lie in a differential module
over $K(z)$ generated by at most $2^{W_\infty}$ functions.
Each individual $G_r$ satisfies a nonzero scalar differential equation
over $K(z)$ of order at most $2^{W_\infty}$.
\end{corollary}
\begin{proof}
Iterate $\mathcal J_0$. There are $2^{w-1}$ positive compositions of
weight $w$, so the empty word and all weights up to $W_\infty$ give
$2^{W_\infty}$ generators, with differentiation preserving their span.
The first $2^{W_\infty}+1$ derivatives of any one element are linearly
dependent over $K(z)$. The assertion is per element: varying rational
prefactors do not give one common scalar equation for every rank.
\end{proof}

\subsection{A double pole, a dilogarithm and an ordinary endpoint}

For $R(n)=(n+1)^{-2}$ a monic basis change from $1,n$ to $1,n+1$
gives
\[
 D_{2,M}=B_M=M\bigl(H_{M+1}^{(2)}-1\bigr)-(H_{M+1}-1)^2.
\]
The general alphabet theorem predicts a classical dilogarithm before
the explicit summation is performed.
\begin{proposition}[A shifted double-pole evaluation]\label{rha:double}
For $|z|<1$,
\begin{equation}\label{rha:double-generator}
 \sum_{M\ge1}B_Mz^M=
 \frac{(3z-2)\operatorname{Li}_2(z)+(1-z)(3L-L^2)-z}
      {z(1-z)^2},\qquad L=-\log(1-z).
\end{equation}
Its zero value is removable, and its ordinary convergent endpoint is
\begin{equation}\label{rha:double-endpoint}
 \sum_{M\ge1}\frac{B_M}{M(M+1)(M+2)}
       =\frac52\zeta(3)-\frac32\zeta(2)-\frac12.
\end{equation}
\end{proposition}
\begin{proof}
The shifted harmonic generators are
\[
 U=\frac{\operatorname{Li}_2(z)}{z(1-z)}-\frac1{1-z},\qquad
 V=\frac{\operatorname{Li}_2(z)+L^2-2L}{z(1-z)}+\frac1{1-z},
\]
for $H_{M+1}^{(2)}-1$ and $(H_{M+1}-1)^2$, respectively.
Then $\theta U-V$ gives \eqref{rha:double-generator}.
The ensemble formula shows $B_M\ge0$; also
$B_M\le M(\zeta(2)-1)$, so the endpoint converges absolutely.
Using $\frac12\int_0^1 z^{M-1}(1-z)^2dz=1/[M(M+1)(M+2)]$,
Tonelli identifies it with $\frac12\int_0^1 N(z)/z^2dz$, where
$N=(3z-2)\operatorname{Li}_2(z)+(1-z)(3L-L^2)-z$.
Here $N/z\to0$ at zero, $N(1)=\zeta(2)-1$, and
\[
 N'=3\operatorname{Li}_2(z)+L^2-2L+2-\frac{2L}{z}.
\]
Integrate by parts, using $\int_0^1\operatorname{Li}_2(z)dz/z=\zeta(3)$,
$\int_0^1L\,dz/z=\zeta(2)$ and $\int_0^1L^2dz/z=2\zeta(3)$.
The last identity follows by extracting the quadratic logarithmic
coefficient of the beta Gamma product; it is also the height-one
relation $\zeta(2,1)=\zeta(3)$. The remaining cancellation must be
kept together:
\[
 \int_0^1\left(\frac1z-\frac L{z^2}\right)dz
       =\left[\frac Lz-L\right]_0^1=-1.
\]
The result is $\frac12[-\zeta(2)+1+5\zeta(3)-2\zeta(2)-2]$,
which proves \eqref{rha:double-endpoint}.
\end{proof}

These classical low-weight Euler reductions are consequences of the
compiler, without a historical novelty claim. Formal coefficient
sharpness, compressed Hahn norms, endpoint cancellation jets and
spectral Stieltjes derivatives in the original delivery retain their
separate proof-audit scopes. None of the alphabet statements settles
the unrestricted moment-prefix conjecture or Gaussian $S_6$/revised $S_8$.
